\scanpage{138}
\mspage{54}

$\underset{\longleftarrow}{\lim}$), et \inserted{l'accent dans} le $\underline{\operatorname{Hom}}^{!'}$ dans le dernier \unsure{mb|membre} indique qu'on exige de plus que le foncteur en question ait un adjoint \uline{à gauche}.

Soient donc \deleted{$M,N$}

\[
M=F(\unsure{X|\mathcal X}),\quad N=F(\unsure{Y|\mathcal Y})
\tag*{(4.1.6)}
\]

et soit $\unsure{Z|\mathcal Z}$ un troisième ps. topos, décrit par

\[
P=F(Z),
\tag*{(4.1.7)}
\]

On veut expliciter, à équivalence près, la catégorie produit $\underline{\operatorname{Hom}}_{\mathrm{ps\,top}}(Z,X)\times\underline{\operatorname{Hom}}_{\mathrm{ps\,top}}(Z,Y)$, i.e.

\[
\underline{\operatorname{Hom}}_{!}^{\deleted{\unsure{!}}}(M,P)^{\mathrm{o}}
\times\operatorname{Hom}_{!}(N,Q)^{\mathrm{o}}
\tag*{(4.1.8)}
\]

\deleted{Si $(f^{*},g^{*})$ est un objet de cette catégorie, je \unsure{dois|désire} leur associer un bifoncteur}

\deleted{\[
\deleted{\unsure{f\times g}}\,f^{*}\boxtimes g^{*}:M\times N\longrightarrow P
\tag*{(4.1.9)}
\]}

\deleted{par}

\deleted{\[
(f^{*}\boxtimes g^{*})(x,y)=
\underbrace{\unsure{g^{*}|f^{*}}(\unsure{a|x})\times f^{\unsure{*|\alpha}}(y)}_{\substack{\text{produit cartésien}\\\text{dans }P}}
\tag*{(4.1.10)}
\]}

Je vais introduire pour cela une catégorie
% archive sheet number: 137
